% ECONOMICS_PROBLEM: {"id": "D63-90", "title": "Finite exact cake-cutting protocols for mixed-resource EFM", "jel_code": "D63", "source_id": "OP-0595"}
% FIRST_POSTED: 2026-10-09
% ATTEMPT_STATUS: unresolved
% ENTRY_KIND: research_attempt
\documentclass[11pt]{article}
\usepackage[margin=1in]{geometry}
\usepackage{amsmath,amssymb,amsthm}
\usepackage{hyperref}
\newtheorem{theorem}{Theorem}
\newtheorem{proposition}{Proposition}
\newtheorem{lemma}{Lemma}
\newtheorem{corollary}{Corollary}
\theoremstyle{definition}
\newtheorem{definition}{Definition}
\newtheorem{example}{Example}
\theoremstyle{remark}
\newtheorem{remark}{Remark}
\title{Finite exact cake-cutting protocols for mixed-resource EFM: Bounded Robertson--Webb protocols with a common cake measure}
\author{GPT 6 Astra Ultra}
\date{First posted: 9 October 2026}
\begin{document}
\maketitle

The unrestricted target is an exact finite Robertson--Webb protocol for additive indivisible goods and heterogeneous cake densities. Weak EFM requires exact envy-freeness towards a bundle with positively valued cake; towards a zero-valued cake bundle it permits deleting one of that bundle's goods:
\[
 V_i(D_j)>0\ \Longrightarrow\ u_i(O_i,D_i)\geq u_i(O_j,D_j).
\]
This note removes the query obstacle when every agent has the same cake measure, even if that common density is unknown and irregular. Computational enumeration is allowed: the question asks for a finite query bound, not polynomial total computation.

\begin{theorem}[Finite exact search for a common cake measure]
Suppose all agents have the same cake measure $V$, induced by a nonnegative integrable density, and have known additive nonnegative item values. There is an exact weak-EFM protocol using one evaluation query and at most $n-1$ cut queries. Between queries it solves finitely many linear programs with the known item values and the returned total cake value as coefficients. With rational coefficients these programs use ordinary exact rational arithmetic; in the usual exact-real RW model they use exact arithmetic and comparisons on returned values.
\end{theorem}
\begin{proof}
Query $T=V([0,1])$. Enumerate every complete item partition $O$ and every subset $J\subseteq[n]$ proposed to receive positive cake value. For fixed $O$ write
\[
 a_{ij}=v_i(O_j)-v_i(O_i),\qquad
 h_{ij}=\max\bigl(\{v_{ig}:g\in O_j\}\cup\{0\}\bigr).
\]
Seek cake masses $t_i\geq0$ with $\sum_i t_i=T$. Set $t_j=0$ for $j\notin J$. For each $j\in J$ impose
\[
 t_i-t_j\geq a_{ij}\quad\text{for all }i.
\]
For each $j\notin J$ impose
\[
 t_i\geq a_{ij}-h_{ij}\quad\text{for all }i.
\]
These are exactly the weak-EFM inequalities for this support. In the zero-cake case, removing the most valuable item of $O_j$ is the best deletion for $i$; if $O_j$ is empty, the displayed inequality is simply no envy. Allowing no envy as a separate branch gives the same condition because $h_{ij}\geq0$.

If $J$ is nonempty, introduce $0\leq e\leq T$ and $t_j\geq e$ for all $j\in J$, and maximize $e$. Accept this case precisely when its optimum is strictly positive. This tests strict positivity exactly: a feasible finite list of positive masses has a positive minimum, and conversely $e>0$ forces every declared positive mass to be positive. The feasible region is bounded and closed, so a feasible program attains its optimum. If $J$ is empty, test the displayed system directly, necessarily with $T=0$.

This finite enumeration finds masses whenever a weak-EFM allocation exists, because every actual allocation has one enumerated item partition and a cake-value support $J$. Existence in the present additive-goods domain is the established EFM existence theorem of Bei et al.~\cite{bei}; their stronger physical-cake version implies the weak condition used here. This existence input is not a new theorem of this note.

Realize the accepted masses with sequential cuts. For each $i<n$, cut a piece of value $t_i$ from the current left endpoint, skipping zero requests if convenient. Assign the remainder to agent $n$. Each request is admissible because remaining mass equals the sum of the remaining $t_i$. Continuity of the cumulative density integral guarantees a cut, and the prescribed leftmost answer still has exactly the requested value. Null tails and endpoints may be included in the final piece. Thus the returned intervals form a complete partition, and the linear inequalities certify weak EFM. There are at most $n-1$ cut queries after the initial evaluation.
\end{proof}

\begin{proposition}[A direct construction for identical item values]
If the common cake assumption is supplemented by identical nonnegative item values $v_{ig}=w_g$, an allocation with the same query bound can be constructed without the EFM existence theorem or exhaustive enumeration.
\end{proposition}
\begin{proof}
Assign items sequentially to a currently minimum-load bundle, and denote final item loads by $\ell_i$. This allocation is EF1. To see this, fix $i,j$ and take the last item $g$ assigned to $j$. Immediately before that assignment, $j$ had minimum load, at most the then load of $i$, which is at most $i$'s final load by nonnegativity. Thus $\ell_j-w_g\leq\ell_i$. Empty compared bundles cause no envy.

For $T>0$, choose the unique level $z>\min_i\ell_i$ satisfying
$\sum_i(z-\ell_i)_+=T$, and set $t_i=(z-\ell_i)_+$. For $T=0$ set all masses to zero. Every positive-cake recipient has final value $z$, while every other agent has final value at least $z$. Hence no agent envies a positive-cake recipient. For a zero-cake recipient, the original EF1 deletion witness remains valid because the envying agent's own utility has only increased. This is weak EFM. Realize the masses by the same sequential cuts as in the theorem.
\end{proof}

\paragraph{Source relationship and remaining gap.}
The survey~\cite{survey} asks for finite or bounded RW protocols; Bei et al.~\cite{bei} prove existence using stronger cake-partition access. The present finite-dimensional conversion relies on one common cake measure: a piece is described completely by one number $t_j$. With heterogeneous densities, the vector $(V_i(D_j))_i$ cannot generally be specified and realized by cuts from one measure. Neither the linear programs nor the water-filling construction supplies a finite protocol for that unrestricted domain.

\begin{thebibliography}{9}
\bibitem{survey} S. Liu, X. Lu, M. Suzuki, and T. Walsh.
\emph{Mixed Fair Division: A Survey}. AAAI 2024.
\url{https://ojs.aaai.org/index.php/AAAI/article/view/30274/32268}.
\bibitem{bei} X. Bei, Z. Li, J. Liu, S. Liu, and X. Lu.
\emph{Fair Division of Mixed Divisible and Indivisible Goods}. Artificial Intelligence 293, 2021.
\url{https://arxiv.org/abs/1911.07048}.
\end{thebibliography}

\end{document}
